\documentclass{article}
\begin{document}
\section*{Pitagorov izrek}
V pravokotnem trikotniku s katetama \(a\) in \(b\) ter hipotenuzo \(c\) velja
\[ a^2 + b^2 = c^2. \]
Hipotenuzo \(c\) lahko izrazimo tudi kot
\[ c^2 = a^2 + b^2 \quad \Longrightarrow \quad c = \sqrt{a^2 + b^2}. \]
\end{document}